\documentclass[11pt,a4paper]{article}

\usepackage[T1]{fontenc}
\usepackage[utf8]{inputenc}
\usepackage{lmodern}

\usepackage[a4paper,margin=1in]{geometry}
\usepackage{setspace}
\setstretch{1.12}

\setlength{\parindent}{15pt}
\setlength{\parskip}{0pt}

\usepackage{lmodern}
\usepackage{amsmath,amssymb,amsfonts}
\usepackage{booktabs}
\usepackage{longtable}
\usepackage{array}
\usepackage{tabularx}
\usepackage{enumitem}
\usepackage{hyperref}
\usepackage{setspace}
\usepackage{ragged2e}

\setlist[itemize]{topsep=3pt,itemsep=2pt,parsep=0pt}
\setlist[enumerate]{topsep=3pt,itemsep=2pt,parsep=0pt}


\title{Glossary of Core Terms, Notation, Interpretation Rules\\
	for the General Theory of Cognitive Structuring Series}
	
\author{
	Kostiantyn Osmolovskyi\\
	\small Independent Researcher, Ukraine\\
	\small ORCID: 0009-0006-3144-7237\\
	\small \texttt{constantinosmol@gmail.com}
}

\date{}

% ---------------------------------------------------------------------------
% Part of a series: General Theory of Cognitive Structuring (GTCS)
% Autor: Kostiantyn Osmolovskyi (ORCID: 0009-0006-3144-7237)
% License: Creative Commons Attribution 4.0 International (CC BY 4.0)
% Repository: https://github.com/k-osmolovskyi/k-osmolovskyi.github.io
% DOI: https://doi.org/10.5281/zenodo.19689204
% ---------------------------------------------------------------------------


\begin{document}
	\maketitle

\begin{center}
	\small
	Technical Note No. 2026-2\\
	Version 1.0
\end{center}

\vspace{0.5em}
	
	\maketitle
	
	\section{Purpose of this document}
	
	This document provides an auxiliary glossary and interpretive guide for the \textit{General Theory of Cognitive Structuring} series. Its purpose is not to replace the formal definitions given in individual papers, but to support accurate external reading, formal reconstruction, and minimal verification of the framework.
	
	The document is written for readers who need a stable interpretive layer across the series, including reviewers, editors, formal verifiers, adjacent-field researchers, and readers attempting cross-paper reconstruction of the architecture.
	
	Its central aim is to reduce category errors and accidental substitutions in the interpretation of key terms. In particular, it is designed to prevent confusion between:
	\begin{itemize}[leftmargin=2em]
		\item global structural conditions and locally accessible signals,
		\item accessibility and representation,
		\item representation and manifestation,
		\item structural admissibility and mere occurrence,
		\item trajectory-dependent state variables and ordinary-language notions of memory,
		\item architectural terms and folk-psychological vocabulary.
	\end{itemize}
	
	Unless explicitly specified otherwise in a given paper, the key terms of the series should be read in an operational, architectural, and regulatory sense.
	
	\section{General interpretive rule}
	
	The central terms of the series do not primarily designate substances, introspectively self-given entities, mental essences, or ordinary-language psychological categories. They designate one or more of the following:
	\begin{itemize}[leftmargin=2em]
		\item structural constraints,
		\item admissibility conditions,
		\item regulatory variables,
		\item accessibility filters,
		\item state components,
		\item stability relations,
		\item trajectory-dependent traces,
		\item manifestation modes of an underlying organization.
	\end{itemize}
	
	A term should therefore be interpreted by its functional role within the architecture rather than by its similarity to everyday language.
	
	\section{Core reading scaffold}
	
	A minimal safe reading scaffold for the series is:
	\[
	\text{global architecture} \rightarrow \text{admissibility/access} \rightarrow \text{regulatory representation} \rightarrow \text{manifestation}.
	\]
	
	This chain is not the full theory, but it provides a reliable interpretive sequence for verification.
	
	In particular, readers should maintain the distinction between four levels:
	\begin{enumerate}[leftmargin=2em]
		\item \textbf{Global structural condition:} the current organization of the system relative to its admissible domain;
		\item \textbf{Accessibility condition:} what part of that organization may become regulatorily available;
		\item \textbf{Regulatory representation:} the internal form in which accessible structural pressure becomes usable for regulation;
		\item \textbf{Manifestation:} the way the current organization becomes expressed, observable, or experience-near.
	\end{enumerate}
	
	A large share of misinterpretations arises from collapsing these levels.
	
	\section{Common misreadings to avoid}
	
	The following substitutions are incorrect unless explicitly argued otherwise in a given text:
	\begin{itemize}[leftmargin=2em]
		\item \textbf{admissibility} $\neq$ truth, correctness, rationality, or mere possibility;
		\item \textbf{coherence} $\neq$ semantic consistency, narrative unity, or subjective harmony;
		\item \textbf{coherence representation} $\neq$ full self-knowledge of global coherence;
		\item \textbf{pre-symbolic admissibility} $\neq$ implicit symbolic content;
		\item \textbf{manifestation} $\neq$ the underlying architecture itself;
		\item \textbf{invariant} $\neq$ necessarily an explicit belief or verbal rule;
		\item \textbf{overload memory} $\neq$ autobiographical memory or semantic memory;
		\item \textbf{regulation} $\neq$ introspection;
		\item \textbf{accessibility} $\neq$ conscious reportability;
		\item \textbf{mode} or \textbf{regime} $\neq$ mood or ordinary-language state label;
		\item \textbf{structural tension} $\neq$ subjective stress by default;
		\item \textbf{compression} $\neq$ deletion or forgetting;
		\item \textbf{level} $\neq$ value hierarchy or depth of insight.
	\end{itemize}
	
	\section{Core glossary}
	
	\subsection{Admissibility}
	
	\textbf{Definition.} Admissibility is the condition under which a process, discrepancy, transformation, representation, or update is structurally allowable within the current architectural organization of the system.
	
	It indicates not merely that something exists or occurs, but that it may enter effective regulation or transformation without violating the relevant stability conditions beyond the operative threshold.
	
	\textbf{Functional role.} Admissibility is a constraint concept. It governs whether some item may participate in regulation, updating, or structurally consequential processing.
	
	\textbf{Does not mean.} Truth, correctness, rationality, social permission, detectability, or mere causal occurrence.
	
	\textbf{Minimal verification cue.} The term is being read correctly if it is tied to structural allowance, thresholds, constraints, stability, and regulatory entry.
	
	\textbf{Short gloss.} A condition of structural allowance under current architectural constraints.
	
	\subsection{Structural Updating}
	
	\textbf{Definition.} Structural updating is a reorganization of the stabilized architecture of processing rather than a local modification of content within an unchanged organization.
	
	The relevant distinction is between change \emph{within} a structure and change \emph{of} the structure.
	
	\textbf{Functional role.} It isolates a level of transformation that is not reducible to ordinary learning, content accumulation, or local adjustment.
	
	\textbf{Does not mean.} Learning in general, memory formation, content acquisition, local adaptation, or parameter change within a preserved architecture.
	
	\textbf{Minimal verification cue.} A correct reading preserves the distinction between content change and structural reorganization.
	
	\textbf{Short gloss.} Reorganization of the stabilized structure of processing, not mere content change within it.
	
	\subsection{Coherence}
	
	\textbf{Definition.} Coherence is a structural relation between the system's current configuration and its admissible region of stable organization.
	
	In the formal line of the series, coherence is not primarily semantic or narrative. It is defined architecturally and, in the formal papers, geometrically.
	
	\textbf{Functional role.} It indexes how the current organization relates to the admissible domain determined by the system's invariants.
	
	\textbf{Does not mean.} Logical consistency of propositions, semantic coherence of discourse, felt unity, or ordinary harmony.
	
	\textbf{Minimal verification cue.} The term is read correctly when tied to admissible region, invariants, stability geometry, or structural distance.
	
	\textbf{Short gloss.} The structural relation between current configuration and the admissible stability domain.
	
	\subsection{Coherence Evaluation}
	
	\textbf{Definition.} Coherence evaluation is the operation or formal determination by which the current configuration is assessed relative to the admissible region of organization.
	
	\textbf{Functional role.} It provides an architectural basis for assessing whether the system remains within, near, or beyond admissible organization.
	
	\textbf{Does not mean.} Self-reflection, verbal judgment, conscious assessment, or semantic interpretation of experience.
	
	\textbf{Minimal verification cue.} Look for evaluation relative to admissible region, structural distance, invariants, or stability geometry.
	
	\textbf{Short gloss.} Architectural assessment of current configuration relative to admissible organization.
	
	\subsection{Coherence Representation}
	
	\textbf{Definition.} Coherence representation is an internal regulatory variable constructed from signals that are available to the system under its accessibility conditions. It is not identical to the full global coherence state.
	
	The system need not regulate by directly accessing the entire structural geometry. It may regulate through an internally available proxy.
	
	\textbf{Functional role.} It links structural condition to actual regulation by providing a usable internal variable.
	
	\textbf{Does not mean.} Direct access to total coherence, full self-model, conscious self-description, or symbolic report by default.
	
	\textbf{Minimal verification cue.} Correct reading preserves the distinction between coherence as structural relation and coherence representation as accessible regulatory proxy.
	
	\textbf{Short gloss.} An internally usable regulatory proxy of coherence, not coherence itself.
	
	\subsection{Pre-Symbolic Admissibility}
	
	\textbf{Definition.} Pre-symbolic admissibility is the condition that determines which discrepancies, pressures, or signals may become regulatorily accessible prior to symbolic articulation, explicit representation, or reportable conceptual form.
	
	It is a gating condition rather than a content category.
	
	\textbf{Functional role.} It restricts the subset of total discrepancies that may enter regulatory processing before symbolic form.
	
	\textbf{Does not mean.} The unconscious in a loose sense, implicit beliefs, pre-linguistic content in general, or nonverbal processing as such.
	
	\textbf{Minimal verification cue.} The term is being read correctly when treated as an access filter rather than as a type of content.
	
	\textbf{Short gloss.} A pre-symbolic gating condition on which discrepancies can become regulatorily available.
	
	\subsection{Admissible Discrepancy Domain}
	
	\textbf{Definition.} The admissible discrepancy domain is the set of discrepancies that pass the relevant admissibility filter and thus become available for regulation.
	
	Not every discrepancy present in the architecture belongs to this domain.
	
	\textbf{Functional role.} It marks the subset of discrepancy space that is effectively available for regulatory engagement.
	
	\textbf{Does not mean.} All discrepancies, all felt discrepancies, all reportable discrepancies, or all symbolized discrepancies.
	
	\textbf{Minimal verification cue.} A correct reading distinguishes discrepancy as such from discrepancy that is admissible for regulation.
	
	\textbf{Short gloss.} The subset of discrepancies that become regulatorily available under admissibility constraints.
	
	\subsection{Restricted Accessibility of Coherence}
	
	\textbf{Definition.} Restricted accessibility of coherence is the principle that global incoherence or structural mismatch does not automatically become accessible to regulation.
	
	The system may be globally outside optimal coherence conditions while lacking regulatorily usable access to that fact.
	
	\textbf{Functional role.} This principle protects the distinction between global structure and accessible signal.
	
	\textbf{Does not mean.} That incoherence does not exist, that the system lacks regulation altogether, or that it is totally blind to its organization.
	
	\textbf{Minimal verification cue.} The term is used correctly when the distinction between global condition, accessible discrepancy, and regulatory signal is preserved.
	
	\textbf{Short gloss.} Global incoherence need not become regulatorily accessible.
	
	\subsection{Representation}
	
	\textbf{Definition.} Representation in this series is generally used in a restricted regulatory sense: an internal variable, form, or organization through which some system-relevant condition becomes usable in processing or regulation.
	
	It should not be assumed to mean symbolic proposition unless explicitly specified.
	
	\textbf{Functional role.} Representation mediates between accessible structure and regulatory or processing consequences.
	
	\textbf{Does not mean.} Necessarily linguistic form, conscious content, symbolic proposition, or full internal picture of reality.
	
	\textbf{Minimal verification cue.} Check whether the represented item is globally present, access-filtered, or already regulatorily usable.
	
	\textbf{Short gloss.} An internal form through which some system-relevant condition becomes usable for regulation or processing.
	
	\subsection{Manifestation}
	
	\textbf{Definition.} Manifestation is the mode in which an underlying architectural and regulatory organization becomes expressed, observable, behaviorally evident, or experience-near.
	
	It is not a separate entity and not an independent causal principle. It is a mode of expression of an underlying organization.
	
	\textbf{Functional role.} Manifestation connects architecture to appearance, expression, and observational layer.
	
	\textbf{Does not mean.} The architecture itself, a personality trait, a symptom by default, behavior alone, or a mental essence.
	
	\textbf{Minimal verification cue.} The term is being used correctly when it refers to expression of organization rather than the organization itself.
	
	\textbf{Short gloss.} The mode in which an underlying architectural-regulatory state becomes expressed.
	
	\subsection{Mode}
	
	\textbf{Definition.} A mode is a relatively stable pattern of system operation characterized by a particular organization of access, regulation, sensitivity, and manifestation.
	
	\textbf{Functional role.} Modes distinguish patterned operational states of the system.
	
	\textbf{Does not mean.} Mood, emotion, or ordinary-language state label.
	
	\textbf{Minimal verification cue.} A mode should alter one or more of accessibility, thresholds, regulatory patterning, or manifestation profile.
	
	\textbf{Short gloss.} A relatively stable pattern of system operation with characteristic access and regulation properties.
	
	\subsection{Regime}
	
	\textbf{Definition.} A regime is a structured regulatory organization governing how the system processes discrepancy, pressure, or representational input under a given configuration.
	
	Where ``mode'' may refer more broadly to operational pattern, ``regime'' often emphasizes the active organization of regulation.
	
	\textbf{Functional role.} It specifies how regulation proceeds under particular structural conditions.
	
	\textbf{Does not mean.} A superficial state label or an ordinary-language condition tag.
	
	\textbf{Minimal verification cue.} Look for relation to regulatory policy, structured processing pattern, or conditional organization of response.
	
	\textbf{Short gloss.} An organized pattern of regulation under a given structural condition.
	
	\subsection{Invariant}
	
	\textbf{Definition.} An invariant is a stabilized structural constraint or organizing condition that participates in defining admissible system organization.
	
	Invariants need not be explicit, verbal, or consciously represented.
	
	\textbf{Functional role.} They define admissible domains, participate in conflict geometry, and anchor organization.
	
	\textbf{Does not mean.} Necessarily a stated belief, verbal rule, metaphysical essence, or permanently fixed feature.
	
	\textbf{Minimal verification cue.} A term is functioning as an invariant if it constrains admissible organization, contributes to stability geometry, or enters structural conflict.
	
	\textbf{Short gloss.} A stabilized structural constraint contributing to admissible organization.
	
	\subsection{Invariant Structure}
	
	\textbf{Definition.} Invariant structure is the organized set or geometry of interacting invariants that defines the system's admissible space.
	
	\textbf{Functional role.} It provides the architecture through which coherence, conflict, and admissibility become formally tractable.
	
	\textbf{Does not mean.} A list of conscious beliefs, a narrative identity, or a static semantic self-description.
	
	\textbf{Minimal verification cue.} Correct reading treats invariants as interacting structural constraints rather than isolated symbolic commitments.
	
	\textbf{Short gloss.} The organized system of interacting invariants that defines admissible organization.
	
	\subsection{Structural Conflict}
	
	\textbf{Definition.} Structural conflict is incompatibility among invariants or structural demands within the architecture.
	
	It is not primarily semantic disagreement but architectural incompatibility.
	
	\textbf{Functional role.} It contributes to structural tension and constrains admissibility.
	
	\textbf{Does not mean.} Ordinary interpersonal conflict, contradiction in discourse alone, or conscious ambivalence by default.
	
	\textbf{Minimal verification cue.} Conflict is structural when defined through incompatibility of organizing constraints.
	
	\textbf{Short gloss.} Architectural incompatibility among structural constraints or invariants.
	
	\subsection{Structural Tension}
	
	\textbf{Definition.} Structural tension is aggregated structural pressure arising from incompatibility among invariants and from the coordination cost required to sustain organization under those incompatibilities.
	
	\textbf{Functional role.} It is a central load variable affecting admissibility, overload, and regulation.
	
	\textbf{Does not mean.} Subjective stress, anxiety, or arousal by default.
	
	\textbf{Minimal verification cue.} The term should be tied to incompatibility, coordination cost, and structural pressure.
	
	\textbf{Short gloss.} Aggregated structural pressure produced by incompatibility and coordination demands.
	
	\subsection{Coordination Cost}
	
	\textbf{Definition.} Coordination cost is the structural cost incurred by the system in maintaining organized processing across potentially competing or nontrivially related constraints.
	
	\textbf{Functional role.} It contributes to total structural tension.
	
	\textbf{Does not mean.} Energetic cost only, social coordination, or narrow computational cost unless explicitly specified.
	
	\textbf{Minimal verification cue.} The term refers to the burden of preserving coordinated organization under structural demand.
	
	\textbf{Short gloss.} The structural cost of maintaining coordinated organization under constraint.
	
	\subsection{Compensability Threshold}
	
	\textbf{Definition.} The compensability threshold is the level of tension the system can absorb without producing a persistent overload trace.
	
	\textbf{Functional role.} It separates compensable pressure from pressure that leaves a trajectory-dependent residual effect.
	
	\textbf{Does not mean.} Tolerance in a colloquial sense, resilience self-report, or moral endurance.
	
	\textbf{Minimal verification cue.} The threshold marks the boundary between absorbable structural pressure and pressure that generates retained overload.
	
	\textbf{Short gloss.} The maximum structural pressure absorbable without persistent overload trace.
	
	\subsection{Overload}
	
	\textbf{Definition.} Overload is the excess of structural tension beyond the system's compensability threshold.
	
	It is not simply high load; it is pressure beyond what can be absorbed without retained consequence.
	
	\textbf{Functional role.} It introduces nontrivial temporal dependence into regulation.
	
	\textbf{Does not mean.} Busyness, tiredness, high activation, or ordinary stress in loose usage.
	
	\textbf{Minimal verification cue.} Correct reading distinguishes tension in general from tension beyond compensable limits.
	
	\textbf{Short gloss.} Structural pressure that exceeds the system's compensable threshold.
	
	\subsection{Overload Memory}
	
	\textbf{Definition.} Overload memory is the trajectory-dependent retention of prior overload. It is a state component that cannot be reduced to current tension alone.
	
	\textbf{Functional role.} It introduces historical dependence into admissibility and regulation.
	
	\textbf{Does not mean.} Autobiographical memory, episodic memory, declarative memory, or semantic memory.
	
	\textbf{Minimal verification cue.} The term is read correctly when treated as a retained structural trace of prior excess rather than as memory content.
	
	\textbf{Short gloss.} Trajectory-dependent retention of prior overload, irreducible to current tension alone.
	
	\subsection{Trajectory Dependence}
	
	\textbf{Definition.} Trajectory dependence means that the current regulatory or admissibility state depends not only on present values but also on the path by which the system arrived there.
	
	\textbf{Functional role.} It prevents collapse of the theory into instantaneous-state description.
	
	\textbf{Does not mean.} Chronological sequence alone or historical narrative in a loose sense.
	
	\textbf{Minimal verification cue.} Two systems with similar present tension may differ due to prior load trajectories.
	
	\textbf{Short gloss.} Dependence of current state on prior path, not present values alone.
	
	\subsection{Stability Region / Admissible Region}
	
	\textbf{Definition.} The stability region, or admissible region, is the set of configurations under which the organization remains structurally admissible according to the relevant criteria.
	
	\textbf{Functional role.} It provides the geometric domain relative to which coherence and admissibility are evaluated.
	
	\textbf{Does not mean.} Comfort zone, habit range, or semantic normality.
	
	\textbf{Minimal verification cue.} Correct reading ties the term to formal domain, structural organization, and admissibility criteria.
	
	\textbf{Short gloss.} The set of configurations under which the system remains structurally admissible.
	
	\subsection{Critical Surface / Critical Boundary}
	
	\textbf{Definition.} The critical surface or critical boundary is the threshold relation separating admissible from non-admissible regions in the relevant state space.
	
	\textbf{Functional role.} It formalizes the transition line beyond which current organization becomes structurally unstable or non-admissible.
	
	\textbf{Does not mean.} A dramatic event, subjective crisis, or phenomenological tipping point by default.
	
	\textbf{Minimal verification cue.} This term belongs to formal threshold geometry.
	
	\textbf{Short gloss.} The threshold boundary separating admissible from non-admissible organization.
	
	\subsection{Configuration}
	
	\textbf{Definition.} A configuration is the current organized state of the system relative to the variables relevant for the theoretical level under discussion.
	
	\textbf{Functional role.} It is the object whose coherence, admissibility, and regulatory implications are evaluated.
	
	\textbf{Does not mean.} Raw sensory input, symbolic self-description, or one single empirical realization only.
	
	\textbf{Minimal verification cue.} Correct reading depends on the abstraction level fixed in the text.
	
	\textbf{Short gloss.} The current organized state of the system at the relevant level of description.
	
	\subsection{Architecture}
	
	\textbf{Definition.} Architecture is the structured organization of constraints, accessibility conditions, regulatory relations, and state variables that define how the system can operate and change.
	
	\textbf{Functional role.} Architecture is the level at which the theory locates admissibility, coherence, overload, and representational constraints.
	
	\textbf{Does not mean.} Anatomy only, implementation substrate only, software architecture in the narrow engineering sense, or phenomenology alone.
	
	\textbf{Minimal verification cue.} Architecture organizes possible states, constraints, access, regulation, and transformation.
	
	\textbf{Short gloss.} The organized system of constraints and regulatory relations that defines possible operation and change.
	
	\subsection{Regulation}
	
	\textbf{Definition.} Regulation is the set of processes by which the system modulates, maintains, redirects, or reorganizes its operation relative to structural pressure and admissibility conditions.
	
	\textbf{Functional role.} It links accessible state information to adjustment or maintenance of organization.
	
	\textbf{Does not mean.} Conscious self-control only, executive function only, or introspection.
	
	\textbf{Minimal verification cue.} Regulation may occur without full symbolic access.
	
	\textbf{Short gloss.} Processes by which the system modulates its operation under structural constraints.
	
	\subsection{Regulatory Accessibility}
	
	\textbf{Definition.} Regulatory accessibility is the condition under which some discrepancy, signal, or variable can enter actual regulation.
	
	\textbf{Functional role.} It distinguishes global presence from regulatorily usable availability.
	
	\textbf{Does not mean.} Conscious awareness, verbal reportability, or salience alone.
	
	\textbf{Minimal verification cue.} A signal may exist without being regulatorily accessible.
	
	\textbf{Short gloss.} Availability of a signal or discrepancy for effective regulation.
	
	\subsection{Expressed Regulatory Representation}
	
	\textbf{Definition.} An expressed regulatory representation is the internally available form in which structurally relevant pressure becomes represented for use in regulation.
	
	This is the operational role often played by variables such as coherence representation.
	
	\textbf{Functional role.} It mediates between admissible structure and regulatory consequence.
	
	\textbf{Does not mean.} Narrative self-understanding, verbal self-report, or full structural transparency.
	
	\textbf{Minimal verification cue.} Its key feature is usability for regulation, not descriptive completeness.
	
	\textbf{Short gloss.} An internal representation usable for regulation of structural pressure.
	
	\subsection{Compression}
	
	\textbf{Definition.} Compression is a structural operation in which multiple invariants or organizing constraints are reorganized into a new, more compact structural unit without simple elimination of their regulatory role.
	
	\textbf{Functional role.} Compression enables reorganization of invariant structure and may alter the geometry of admissible organization.
	
	\textbf{Does not mean.} Deletion, forgetting, or data compression in a narrowly technical sense.
	
	\textbf{Minimal verification cue.} Correct reading preserves reorganization rather than mere removal.
	
	\textbf{Short gloss.} Reorganization of multiple structural constraints into a new compact unit.
	
	\subsection{Level}
	
	\textbf{Definition.} A level is a class of architectural organization defined not by content amount or subjective depth but by the form of admissible stability geometry.
	
	\textbf{Functional role.} It provides a formal criterion for distinguishing structurally non-equivalent organizations.
	
	\textbf{Does not mean.} Importance, hierarchy of value, or introspective depth.
	
	\textbf{Minimal verification cue.} A level must be tied to equivalence or non-equivalence of structural admissibility geometry.
	
	\textbf{Short gloss.} A class of organization defined by the form of its admissible stability geometry.
	
	\subsection{Level Transition}
	
	\textbf{Definition.} A level transition is a non-equivalent transformation of the system's admissible organization such that the resulting stability geometry cannot be treated as merely a reparameterized version of the prior one.
	
	\textbf{Functional role.} It distinguishes true structural reorganization from quantitative variation within one level.
	
	\textbf{Does not mean.} More intensity, more information, or stronger experience.
	
	\textbf{Minimal verification cue.} The criterion is non-equivalence of admissible geometry.
	
	\textbf{Short gloss.} A non-equivalent change in the geometry of admissible organization.
	
	\subsection{Identity as Regulatory Attractor}
	
	\textbf{Definition.} Within this framework, identity is not treated as a metaphysical self or narrative self-description, but as a relatively stable regulatory attractor shaped by invariants, admissibility, and bounded organization.
	
	\textbf{Functional role.} It names a stabilized tendency of organization rather than a substance.
	
	\textbf{Does not mean.} Personal essence, autobiographical story, or explicit self-concept only.
	
	\textbf{Minimal verification cue.} Correct reading treats identity as a regulatory stability phenomenon.
	
	\textbf{Short gloss.} A relatively stable regulatory attractor of organization.
	
	\subsection{Phenomenal or Felt Layer}
	
	\textbf{Definition.} Where used, the phenomenal or felt layer refers not to an autonomous metaphysical realm but to the manner in which some regulatory condition becomes experience-near or phenomenally accessible.
	
	\textbf{Functional role.} It belongs to manifestation and access analysis, not to ontology.
	
	\textbf{Does not mean.} A metaphysical qualia thesis by default, an experiential substance, or direct access to total architecture.
	
	\textbf{Minimal verification cue.} Correct reading treats the phenomenal layer as an architectural effect or manifestation layer.
	
	\textbf{Short gloss.} The experience-near manifestation of regulatory organization, not the architecture itself.
	
	\subsection{Perception}
	
	\textbf{Definition.} Within the broader framework, perception is not reduced to passive input registration. It is understood as a structured result of architecture, admissibility, and regulatory accessibility conditions.
	
	\textbf{Functional role.} Perception belongs to the family of manifestations and representationally mediated system-world relations shaped by architecture.
	
	\textbf{Does not mean.} Raw input alone, transparent access to reality, or purely bottom-up registration.
	
	\textbf{Minimal verification cue.} Perception is treated correctly when constrained by architecture, access, and representation.
	
	\textbf{Short gloss.} Architecturally shaped system-world registration under admissibility and regulatory constraints.
	
	\subsection{Thought}
	
	\textbf{Definition.} Thought is not assumed to be the master layer of the system. It is one possible manifestation and one possible medium of regulatory organization under given access conditions.
	
	\textbf{Functional role.} It belongs to the manifestation and processing layer rather than automatically to the structural layer.
	
	\textbf{Does not mean.} The full architecture, the source of all regulation, or necessarily conscious reasoning.
	
	\textbf{Minimal verification cue.} Thought should not be used as a synonym for organization itself.
	
	\textbf{Short gloss.} A manifestation-level process that may participate in regulation but does not exhaust system architecture.
	
	\subsection{Behavior}
	
	\textbf{Definition.} Behavior is an external manifestation of regulatory organization, not a sufficient description of that organization.
	
	\textbf{Functional role.} It is one output layer through which architecture may become partially observable.
	
	\textbf{Does not mean.} The whole system, direct transparency of organization, or a self-explanatory causal level.
	
	\textbf{Minimal verification cue.} Behavior must be treated as manifestation, not as the architecture itself.
	
	\textbf{Short gloss.} An external manifestation of underlying regulatory organization.
	
	\section{Relational distinctions that must remain explicit}
	
	\subsection{Admissibility vs accessibility}
	
	\begin{itemize}[leftmargin=2em]
		\item \textbf{Admissibility} concerns structural allowance.
		\item \textbf{Accessibility} concerns whether something becomes available to a process or regulatory layer.
	\end{itemize}
	
	Something may be structurally relevant yet inaccessible. Something may be locally accessible without implying structural admissibility for updating.
	
	\subsection{Coherence vs coherence representation}
	
	\begin{itemize}[leftmargin=2em]
		\item \textbf{Coherence} is a structural relation.
		\item \textbf{Coherence representation} is an internal regulatory variable constructed from accessible signals.
	\end{itemize}
	
	The second is not identical to the first.
	
	\subsection{Global state vs admissible discrepancy}
	
	\begin{itemize}[leftmargin=2em]
		\item \textbf{Global state} includes all structurally relevant conditions.
		\item \textbf{Admissible discrepancy} is only the subset that passes the relevant access filter.
	\end{itemize}
	
	\subsection{Representation vs manifestation}
	
	\begin{itemize}[leftmargin=2em]
		\item \textbf{Representation} is an internal usable form.
		\item \textbf{Manifestation} is the way organization becomes expressed or observable.
	\end{itemize}
	
	A representation may contribute to manifestation, but the two are not identical.
	
	\subsection{Tension vs overload}
	
	\begin{itemize}[leftmargin=2em]
		\item \textbf{Tension} is structural pressure.
		\item \textbf{Overload} is pressure beyond the compensability threshold.
	\end{itemize}
	
	Not all tension is overload.
	
	\subsection{Overload vs overload memory}
	
	\begin{itemize}[leftmargin=2em]
		\item \textbf{Overload} is present excess pressure.
		\item \textbf{Overload memory} is retained trace of prior excess.
	\end{itemize}
	
	They are not identical.
	
	\subsection{Learning vs structural updating}
	
	\begin{itemize}[leftmargin=2em]
		\item \textbf{Learning} may occur within existing organization.
		\item \textbf{Structural updating} changes the stabilized organization itself.
	\end{itemize}
	
	\subsection{Invariant vs explicit belief}
	
	\begin{itemize}[leftmargin=2em]
		\item \textbf{Invariant} is a structural constraint.
		\item An \textbf{explicit belief} may or may not express one invariant.
		\item Many invariants need not appear in symbolic form.
	\end{itemize}
	
	\subsection{Global geometry vs local regulation}
	
	\begin{itemize}[leftmargin=2em]
		\item \textbf{Global geometry} concerns the system's position relative to admissible organization.
		\item \textbf{Local regulation} concerns what the system can actually process and modulate through accessible signals.
	\end{itemize}
	
	The latter does not imply full access to the former.
	
	\section{Notation companion}
	
	This section is interpretive rather than exhaustive. It gives readers a stable orientation to frequently recurring notation across the series.
	
	\subsection{$X$}
	
	\textbf{Meaning.} Configuration space or state space at the relevant level of description.
	
	\textbf{Interpretive note.} $X$ does not necessarily mean physical space, neural space, or phenomenological space specifically. It is the structured domain in which current configurations are considered.
	
	\subsection{$x_t$}
	
	\textbf{Meaning.} Current configuration of the system at time $t$.
	
	\textbf{Interpretive note.} $x_t$ is the current organized state relative to the variables fixed by the model, not necessarily a raw empirical observation.
	
	\subsection{$I_t$}
	
	\textbf{Meaning.} The current set of invariants at time $t$.
	
	\textbf{Interpretive note.} These are stabilized structural constraints relevant to admissible organization.
	
	\subsection{$U_X(I_t)$}
	
	\textbf{Meaning.} The admissible or stable region in configuration space $X$ determined by the current invariant structure $I_t$.
	
	\textbf{Interpretive note.} This is the domain relative to which coherence is evaluated. It is not a semantic category but a structural region.
	
	\subsection{$C_I(x_t)$}
	
	\textbf{Meaning.} A coherence measure defined as the structural relation, often formalized as distance, between current configuration $x_t$ and admissible region $U_X(I_t)$:
	\[
	C_I(x_t)=\mathrm{dist}(x_t,U_X(I_t)).
	\]
	
	\textbf{Interpretive note.} This is not narrative or semantic coherence. It is geometric or structural coherence.
	
	\subsection{$a_i(x)$}
	
	\textbf{Meaning.} Local admissibility function for invariant $I_i$, indicating the degree to which configuration $x$ satisfies that invariant.
	
	\textbf{Interpretive note.} These functions support the move from local invariant satisfaction to global structural relations.
	
	\subsection{$c_{ij}(x)$}
	
	\textbf{Meaning.} Local conflict term between invariants $I_i$ and $I_j$ at configuration $x$.
	
	\textbf{Interpretive note.} This expresses incompatibility at the level of structural constraints rather than contradiction in language.
	
	\subsection{$K_t$}
	
	\textbf{Meaning.} Conflict load or incompatibility load at time $t$.
	
	\textbf{Interpretive note.} This is the component of structural pressure arising from incompatibility among invariants or structural demands.
	
	\subsection{$C_t$ in the tension decomposition}
	
	\textbf{Meaning.} Coordination cost at time $t$.
	
	\textbf{Interpretive note.} In contexts where structural tension is decomposed as
	\[
	T_t = K_t + C_t,
	\]
	this $C_t$ denotes coordination cost, not coherence. Context must determine which $C$-notation is in use.
	
	\subsection{$T_t$}
	
	\textbf{Meaning.} Structural tension at time $t$.
	
	A recurring form is:
	\[
	T_t = K_t + C_t.
	\]
	
	\textbf{Interpretive note.} This is an aggregate structural load variable, not subjective stress by default.
	
	\subsection{$\theta(R_t,\bar h_t)$}
	
	\textbf{Meaning.} Compensability threshold as a function of current resource state $R_t$ and rigidity parameter $\bar h_t$.
	
	\textbf{Interpretive note.} It captures how much tension can be absorbed without generating persistent overload trace.
	
	\subsection{$u_t$}
	
	\textbf{Meaning.} Excess overload term at time $t$, often defined as
	\[
	u_t = \max(0, T_t - \theta(R_t,\bar h_t)).
	\]
	
	\textbf{Interpretive note.} This is not total tension but only the excess beyond compensability.
	
	\subsection{$L_t$}
	
	\textbf{Meaning.} Overload memory or retained overload state at time $t$.
	
	\textbf{Interpretive note.} This is a trajectory-dependent state variable, not autobiographical memory content.
	
	\subsection{$A(\cdotp), \quad A_t$}
	
	\textbf{Meaning.} Admissibility operator or admissibility score at time $t$, often written as a function of tension, overload memory, resource, rigidity, and related parameters.
	
	\textbf{Interpretive note.} $A(\cdotp)$ quantifies admissibility conditions under the formal model. Its role is regulatory and structural rather than semantic.
	
	\subsection{$U(S_t)$}
	
	\textbf{Meaning.} Admissible or stability region defined in terms of a regulatory state vector $S_t$, often at the level of aggregated state variables such as tension and overload memory.
	
	\textbf{Interpretive note.} This notation marks the admissible region in a reduced regulatory state space.
	
	\subsection{$T_{\mathrm{crit}}(L;\cdots)$}
	
	\textbf{Meaning.} Critical tension boundary as a function of overload memory and other parameters.
	
	\textbf{Interpretive note.} This defines the threshold separating admissible and non-admissible organization in the reduced state space.
	
	\subsection{$\Pi_t^{ps}$}
	
	\textbf{Meaning.} Pre-symbolic admissibility operator or gate at time $t$.
	
	\textbf{Interpretive note.} This notation marks the filtering condition that determines which discrepancies become regulatorily available before symbolic articulation.
	
	\subsection{$Z_t$}
	
	\textbf{Meaning.} A raw discrepancy or state summary before admissibility filtering, depending on the paper's local notation.
	
	\textbf{Interpretive note.} $Z_t$ generally names a pre-filter representational or discrepancy-level variable.
	
	\subsection{$\mathcal{D}_t$}
	
	\textbf{Meaning.} Admissible discrepancy domain at time $t$.
	
	\textbf{Interpretive note.} This is the subset of discrepancy space that passes the relevant admissibility condition.
	
	\subsection{$\widehat{C}_t$}
	
	\textbf{Meaning.} Expressed coherence representation or coherence-related regulatory proxy at time $t$.
	
	\textbf{Interpretive note.} This variable is especially important to read correctly: it is not the full global coherence geometry, but the internally available representation used in regulation.
	
	\subsection{$\widetilde{C}_t^A$}
	
	\textbf{Meaning.} An extended or transformed regulatory coherence-related variable, depending on the paper's local use.
	
	\textbf{Interpretive note.} When present, it should be read as a derived representational quantity, not as the global structure itself.
	
	\subsection{$F(S_t)$}
	
	\textbf{Meaning.} A mapping from relevant state variables to a regulatory representation, often used in forms such as
	\[
	\widehat{C}_t = F(S_t).
	\]
	
	\textbf{Interpretive note.} $F$ indicates construction of usable internal representation from structurally relevant signals.
	
	\subsection{$\Psi$}
	
	\textbf{Meaning.} A regulatory or organizational mapping from an expressed representation to downstream invariant organization, policy, or updating-related consequence, depending on local formalization.
	
	\textbf{Interpretive note.} Read $\Psi$ functionally from context; it is generally a downstream regulatory map.
	
	\subsection{$\rho$}
	
	\textbf{Meaning.} A regime-selection or policy-selection map in contexts where different regulatory regimes are modeled.
	
	\textbf{Interpretive note.} This notation highlights selection of operational regime from an available internal state or representation.
	
	\subsection{$F(S)$ under compression}
	
	\textbf{Meaning.} Where used in compression contexts, $F(S)$ may denote the result of a compression operation over a subset of invariants.
	
	\textbf{Interpretive note.} Do not confuse this with ordinary summarization. In compression papers, it denotes structural reorganization.
	
	\section{Disallowed substitutions table}
	
	\begin{longtable}{>{\RaggedRight\arraybackslash}p{0.34\textwidth} >{\RaggedRight\arraybackslash}p{0.58\textwidth}}
		\toprule
		\textbf{Series term} & \textbf{Must not be reduced to} \\
		\midrule
		\endfirsthead
		\toprule
		\textbf{Series term} & \textbf{Must not be reduced to} \\
		\midrule
		\endhead
		Admissibility & truth, correctness, awareness, mere possibility \\
		Structural updating & learning, memory formation, content revision \\
		Coherence & harmony, semantic consistency, narrative unity \\
		Coherence representation & full self-knowledge, direct access to global structure \\
		Pre-symbolic admissibility & implicit symbolic content, pre-linguistic content in general \\
		Admissible discrepancy domain & all discrepancies, all felt tensions \\
		Restricted accessibility of coherence & absence of incoherence, total blindness \\
		Manifestation & underlying architecture, behavior alone, essence \\
		Mode & mood, vague state label \\
		Regime & superficial label for condition \\
		Invariant & explicit belief, verbal rule, immutable essence \\
		Structural conflict & interpersonal conflict, contradiction in discourse alone \\
		Structural tension & stress, anxiety, subjective discomfort by default \\
		Coordination cost & energy expenditure alone, social coordination \\
		Compensability threshold & tolerance, willpower, resilience report \\
		Overload & busyness, fatigue, high stimulation \\
		Overload memory & autobiographical memory, semantic memory \\
		Trajectory dependence & chronology alone \\
		Stability region & comfort zone, normality \\
		Critical boundary & crisis experience, dramatic event \\
		Architecture & anatomy only, implementation only \\
		Regulation & introspection, executive control only \\
		Regulatory accessibility & consciousness, verbal reportability \\
		Compression & deletion, forgetting, simplification in the ordinary sense \\
		Level & hierarchy of value, degree of worth \\
		Level transition & stronger intensity, larger quantity \\
		Identity as regulatory attractor & essence, story of self, autobiographical self \\
		Phenomenal layer & metaphysical qualia substance \\
		Perception & raw input registration \\
		Thought & total architecture, master process \\
		Behavior & the whole system \\
		\bottomrule
	\end{longtable}
	
	\section{Minimal verification protocol}
	
	A reader attempting to verify the use of terms in the series may apply the following minimal protocol.
	
	\subsection*{Step 1}
	Identify whether the term refers primarily to:
	\begin{itemize}[leftmargin=2em]
		\item structure,
		\item access,
		\item representation,
		\item regulation,
		\item manifestation,
		\item trajectory-dependent state,
		\item threshold geometry.
	\end{itemize}
	
	\subsection*{Step 2}
	Identify whether the term functions as:
	\begin{itemize}[leftmargin=2em]
		\item a constraint,
		\item a variable,
		\item a state component,
		\item a relation,
		\item a map,
		\item a mode of expression.
	\end{itemize}
	
	\subsection*{Step 3}
	Check whether ordinary-language meaning has replaced architectural meaning.
	
	\subsection*{Step 4}
	Check whether neighboring levels have been collapsed:
	\begin{itemize}[leftmargin=2em]
		\item global state with accessible state,
		\item accessible state with representation,
		\item representation with manifestation,
		\item manifestation with structure,
		\item tension with overload,
		\item overload with overload memory.
	\end{itemize}
	
	\subsection*{Step 5}
	Where notation is present, treat notation as the stronger interpretive anchor and prose as its explanatory support.
	
	\subsection*{Step 6}
	Where a term appears in multiple papers, prefer the narrower functional reading that preserves cross-paper consistency rather than the broader ordinary-language reading.
	
	\section{Priority rules for cross-paper consistency}
	
	When reconstructing the theory across multiple papers, the following priority rules are recommended:
	\begin{enumerate}[leftmargin=2em]
		\item \textbf{Formal role overrides ordinary-language resemblance.} If a term sounds familiar but functions differently in the framework, its formal role takes priority.
		\item \textbf{Cross-paper consistency overrides local rhetorical convenience.} A local phrasing should not be allowed to overturn the stable architectural function of a term across the series.
		\item \textbf{Operator-level distinctions must be preserved.} A discrepancy, an admissibility gate, a representation, and a manifestation are not interchangeable.
		\item \textbf{Trajectory-dependent variables must not be collapsed into instantaneous quantities.} In particular, overload memory must not be reduced to current tension.
		\item \textbf{Manifestational vocabulary must not be mistaken for ontological commitment.} Terms referring to felt, phenomenal, behavioral, or experience-near layers should be read as manifestation-level descriptions unless explicitly formalized otherwise.
	\end{enumerate}
	
	\section{Paper-to-paper mapping table}
	
	This table is intended as a cross-series orientation aid. It indicates where a term or formal component is primarily introduced, clarified, or extended. The table is interpretive rather than bibliographically exhaustive.
	
	\begin{longtable}{>{\RaggedRight\arraybackslash}p{0.25\textwidth} >{\RaggedRight\arraybackslash}p{0.24\textwidth} >{\RaggedRight\arraybackslash}p{0.16\textwidth} >{\RaggedRight\arraybackslash}p{0.25\textwidth}}
		\toprule
		\textbf{Paper} & \textbf{Primary role in the series} & \textbf{Status} & \textbf{Main terms/components} \\
		\midrule
		\endfirsthead
		\toprule
		\textbf{Paper} & \textbf{Primary role in the series} & \textbf{Status} & \textbf{Main terms/components} \\
		\midrule
		\endhead
		\emph{Structural Updating and the Limits of Cognitive Change} & Entry-point distinction isolating the gap between processing access and structural updating & Conceptual & admissibility, structural updating, limit of change, distinction from learning/plasticity \\
		\emph{Coherence Evaluation in Cognitive Systems} & Introduces coherence as a structural-geometric relation to admissible organization & Formal core & coherence, coherence evaluation, configuration, admissible region, $C_I(x_t)$ \\
		\emph{Structural Admissibility in Cognitive Systems} & Provides the main formal architecture for admissibility, thresholds, and stability geometry & Formal core & admissibility operator, stability region, critical boundary, tension, overload memory, level, level transition \\
		\emph{Overload Formation in Cognitive Processing} & Formalizes overload accumulation and non-reducible trajectory dependence & Formal extension & overload, overload memory, excess term $u_t$, retention dynamics, path dependence \\
		\emph{Trajectory-Dependent Regulation in Cognitive Systems} & Extends the framework from state description to trajectory-sensitive regulation & Formal extension & trajectory dependence, regulation across paths, dynamic admissibility consequences \\
		\emph{Identity as a Regulatory Attractor} & Recasts identity as a stability phenomenon within bounded regulation & Formal-conceptual extension & identity as regulatory attractor, bounded organization, persistence regions \\
		\emph{Invariants in Cognitive Architectures} & Formalizes invariants as structural constraints defining admissible organization & Formal core & invariant, invariant structure, local admissibility, conflict geometry \\
		\emph{Structural Compression in Cognitive Architectures} & Introduces compression as structural reorganization of invariant organization & Formal extension & compression, reorganization of invariant sets, level shift through geometric deformation \\
		\emph{Emergence of Coherence Representation} & Introduces internal regulatory proxy of structural coherence & Formal-conceptual bridge & coherence representation, expressed regulatory representation, mapping from structural signals to internal variable \\
		\emph{Coherence Representation in Multi-System Regulation} & Extends coherence representation beyond a single system architecture & Formal extension & transfer of coherence representation, multi-system regulation, inter-architectural usability \\
		\emph{Pre-Symbolic Admissibility in Cognitive Systems} & Introduces the pre-symbolic gate governing regulatory accessibility & Formal core / access layer & pre-symbolic admissibility, admissibility filtering, $\Pi_t^{ps}$, admissible discrepancy domain \\
		\emph{Restricted Accessibility of Coherence in Cognitive Systems} & Establishes that global incoherence does not imply regulatorily accessible incoherence & Proposition-level extension & restricted accessibility, global state vs accessible discrepancy, non-injectivity of access \\
		\emph{Inter-System Conflict Geometry in Cognitive Systems} & Extends conflict analysis to relations among interacting systems & Formal extension & inter-system conflict, cross-system constraint geometry, architectural incompatibility across systems \\
		\emph{General Theory of Cognitive Structuring} & Integrative synthesis of the full architecture and its cross-domain scope & Synthetic / unifying theory & series-wide integration of admissibility, coherence, overload, invariants, compression, representation, levels \\
		\bottomrule
	\end{longtable}
	
	\newpage
	\section{Short reference glossary}
	
	For quick use, the following compact list may serve as a one-page reference.
	
	\begin{itemize}[leftmargin=2em]
		\item \textbf{Admissibility} --- structural allowance under current architectural constraints.
		\item \textbf{Structural updating} --- reorganization of stabilized structure, not mere content change.
		\item \textbf{Coherence} --- structural relation to admissible stability domain.
		\item \textbf{Coherence evaluation} --- assessment of configuration relative to admissible organization.
		\item \textbf{Coherence representation} --- internal regulatory proxy of coherence.
		\item \textbf{Pre-symbolic admissibility} --- pre-symbolic gating of regulatory access.
		\item \textbf{Admissible discrepancy domain} --- discrepancies available for regulation after admissibility filtering.
		\item \textbf{Restricted accessibility of coherence} --- global incoherence may remain regulatorily inaccessible.
		\item \textbf{Manifestation} --- expressed mode of underlying organization.
		\item \textbf{Mode} --- relatively stable operational pattern.
		\item \textbf{Regime} --- organized pattern of regulation.
		\item \textbf{Invariant} --- stabilized structural constraint.
		\item \textbf{Structural conflict} --- incompatibility among structural constraints.
		\item \textbf{Structural tension} --- aggregated structural pressure.
		\item \textbf{Coordination cost} --- burden of preserving organized coordination under constraint.
		\item \textbf{Compensability threshold} --- maximum pressure absorbable without persistent overload trace.
		\item \textbf{Overload} --- pressure beyond compensable threshold.
		\item \textbf{Overload memory} --- retained trace of prior overload.
		\item \textbf{Trajectory dependence} --- dependence on prior path, not only current values.
		\item \textbf{Stability region} --- set of structurally admissible configurations.
		\item \textbf{Critical boundary} --- threshold separating admissible and non-admissible organization.
		\item \textbf{Representation} --- internal usable form for regulation or processing.
		\item \textbf{Regulatory accessibility} --- availability for effective regulation.
		\item \textbf{Compression} --- reorganization of multiple constraints into a new compact unit.
		\item \textbf{Level} --- class of organization defined by admissible geometry.
		\item \textbf{Level transition} --- non-equivalent change in admissible geometry.
		\item \textbf{Identity as regulatory attractor} --- relatively stable regulatory tendency of organization.
		\item \textbf{Phenomenal layer} --- experience-near manifestation of regulatory organization.
		\item \textbf{Perception} --- architecturally shaped system-world registration under admissibility constraints.
		\item \textbf{Thought} --- manifestation-level process participating in regulation.
		\item \textbf{Behavior} --- external manifestation of regulatory organization.
	\end{itemize}
	
	\section{Closing note}
	
	This glossary is intended as a verificatory and interpretive aid. It should be used together with the formal definitions, propositions, and notation provided in the individual papers. Where local formulations differ in emphasis, the preferred reading is the one that preserves architectural consistency across the series.
	
\end{document}